\subsection{Charging compact bags}
\label{hc:section}

We prove the compact-bag estimate, including its geometric charging
step.  The argument follows the construction underlying
\cite[Proposition~7.1 and Section~7.1]{AOG}; the proof below includes the
boundary conventions, the positive-part accounting, and the invariants
needed when the scale changes.

\begin{proposition}[Compact-bag charging]
\label{prop:hierarchy-compact}
Let \(G=(V,E)\) be a finite \(k\)-edge-connected multigraph, and let
\(X_v\in Q=[0,1]^h\) for \(v\in V\).
Suppose \((\mathcal F_i)\) is a finite or countable
\(\lambda\)-geometric sequence of nonempty families of
\(\beta\)-compact bags, of respective types \((\delta_i,\Delta_i)\), where
\(0<\lambda\le 1/12\) and \(\beta\ge36\).  Thus all balls in
\(\mathcal F_i\) are mutually disjoint ordinary balls
\(B_Q(X_v,\delta_i)\), every bag \(D\in\mathcal F_i\) has at least two
balls and center diameter at most \(\Delta_i\), and
\[
 \beta\Delta_i\le |D|\delta_i,
 \qquad
 \Delta_{i+1}\le\lambda\delta_i.
\]
Then
\begin{equation}
\label{hc:conclusion}
 \frac{k}{4}\sum_i\delta_i\sum_{D\in\mathcal F_i}|D|
 \le \sum_{e\in E}\|X_e\|_1,
 \qquad X_{\{u,v\}}=X_u-X_v.
\end{equation}
\end{proposition}

An ordinary ball is \(B_Q(x,r)=\{y\in Q:\|y-x\|_1<r\}\), including
its center.  A hollowed ball is the strict shell
\[
 B_Q(x;a,b)=\{y\in Q:a<\|y-x\|_1<b\},
 \qquad 0\le a<b.
\]
We call \(b-a\) its width.  For numerical formulas an ordinary ball of
radius \(b\) has inner radius \(a=0\) and width \(b\); it remains an
ordinary ball, rather than a shell punctured at its center.  A shell of
zero width is discarded.

\begin{lemma}[Length bound for disjoint hollowed balls]
\label{hc:length}
Let \(G=(V,E)\) be a finite \(k\)-edge-connected multigraph, where
\(k>0\), and let \(X_v\in Q=[0,1]^h\) for \(v\in V\).
Let \(\mathcal Z\) be a finite collection of pairwise disjoint ordinary
or hollowed balls in \(Q\).  For each \(H\in\mathcal Z\), write
\(x_H,a_H,b_H\) for its center and numerical inner and outer radii.
Suppose there are vertices \(u_H,v_H\in V\) such that
\begin{equation}
\label{hc:witness}
 \|X_{u_H}-x_H\|_1\le a_H,
 \qquad \|X_{v_H}-x_H\|_1\ge b_H.
\end{equation}
Then
\begin{equation}
\label{hc:length-bound}
 k\sum_{H\in\mathcal Z}(b_H-a_H)
 \le \sum_{\{u,v\}\in E}\|X_u-X_v\|_1.
\end{equation}
\end{lemma}

\begin{proof}
For \(H\in\mathcal Z\), clamp the radial distance to its radial
interval:
\[
 f_H(y)=\min\{b_H,\max\{a_H,\|y-x_H\|_1\}\}.
\]
For every \(a_H<t<b_H\), the set
\(S_{H,t}=\{v:\|X_v-x_H\|_1<t\}\) is nonempty and proper, by
\eqref{hc:witness}.  Its edge boundary therefore has at least \(k\)
edges, with multiplicity.  Integrating this finite step function gives
\begin{equation}
\label{hc:cut-integral}
 k(b_H-a_H)
 \le \int_{a_H}^{b_H}|E(S_{H,t},V\setminus S_{H,t})|\,dt
 =\sum_{\{u,v\}\in E}|f_H(X_u)-f_H(X_v)|.
\end{equation}
Indeed, the contribution of an edge to the integral is exactly the
length of the interval of thresholds lying between its two endpoint
distances, intersected with \((a_H,b_H)\).

Fix an edge \(e=\{u,v\}\), put \(L_e=\|X_u-X_v\|_1\), and let
\(\gamma(s)=(1-s)X_u+sX_v\), \(0\le s\le1\).  The segment lies in
the convex set \(Q\).  Every \(g_H=f_H\circ\gamma\) is a continuous,
finitely piecewise affine function: each coordinate in its defining
norm is an absolute value of an affine function, and clamping adds
only finitely many breakpoints.  The reverse triangle inequality and
the fact that clamping does not increase distances show that
\(g_H\) is \(L_e\)-Lipschitz.

Take a common finite subdivision \(0=s_0<\cdots<s_m=1\) on which all
the \(g_H\) are affine.  If \(g_H\) is nonconstant on one of these
intervals, its value in the interval's interior lies strictly between
\(a_H\) and \(b_H\).  The corresponding points of \(\gamma\) belong
to \(H\).  Pairwise disjointness permits at most one such nonconstant
function on each interval.  Hence
\begin{align*}
 \sum_{H\in\mathcal Z}|g_H(1)-g_H(0)|
 &\le \sum_{j=1}^m\sum_{H\in\mathcal Z}
       |g_H(s_j)-g_H(s_{j-1})|\\
 &\le \sum_{j=1}^m L_e(s_j-s_{j-1})=L_e.
\end{align*}
This also handles a segment lying along a shell boundary: its clamped
radial function is constant there.  Summing
\eqref{hc:cut-integral} over \(H\) and then applying the last bound to
each edge proves \eqref{hc:length-bound}.
\end{proof}

\begin{proof}[Proof of Proposition~\ref{prop:hierarchy-compact}]
It suffices to prove the claim for each finite initial sequence
\(\mathcal F_1,\ldots,\mathcal F_q\).  For this proof only, set
\(\Delta_{q+1}=0\).  Each family is finite because its positive-radius
balls are disjoint and their centers are drawn from the finite set
\(\{X_v:v\in V\}\).

We first record the scale inequalities.  Two disjoint balls of radius
\(\delta_i\), with centers \(z,z'\in Q\), satisfy
\(\|z-z'\|_1\ge2\delta_i\): otherwise their midpoint, which belongs
to \(Q\), would lie in both balls.  A compact bag contains two such
balls.  Consequently
\begin{equation}
\label{hc:scales}
 2\delta_i\le\Delta_i,
 \qquad
 6\Delta_{i+1}\le\frac{\delta_i}{2},
 \qquad
 |D|\delta_i\ge \beta\Delta_i
 \quad(D\in\mathcal F_i).
\end{equation}

\paragraph{The invariant and the order of processing.}
We construct a finite collection \(\mathcal Z\), initially empty,
and process the families in the order \(1,\ldots,q\).
During phase \(i\), call an original ball inserted from
\(\mathcal F_i\) \emph{fresh}.  Every other live object is called
\emph{old}, including any shell produced by splitting an old object
during this phase.  A fresh ball is kept intact until the phase ends.
Write \([s]_+=\max\{s,0\}\).  Assign capacities
\begin{equation}
\label{hc:capacity}
 p_i(H)=
 \begin{cases}
  \delta_i-6\Delta_{i+1},&H\text{ is fresh},\\
  [w(H)-6\Delta_i]_+,&H\text{ is old},
 \end{cases}
 \qquad P_i(\mathcal Z)=\sum_{H\in\mathcal Z}p_i(H),
\end{equation}
where \(w(H)\) denotes width.  All capacities are nonnegative and
at most the object's width.  At every step we maintain:
\begin{enumerate}[label=(\roman*),leftmargin=*]
\item the live objects are pairwise disjoint;
\item their centers are embedded vertices, and every live object has
      the two witnesses in \eqref{hc:witness};
\item
\begin{equation}
\label{hc:potential-invariant}
 P_i(\mathcal Z)\ge
 \frac14\sum_{D\text{ already processed}}|D|\delta(D),
\end{equation}
where \(\delta(D)\) is the radius of the balls of \(D\).
\end{enumerate}
These assertions hold initially.  We verify them for each operation
and for each change of phase.

\paragraph{Processing a bag with an interior ball.}
For an unprocessed ball \(B_Q(z,\delta_i)\) and a live object \(H\)
with center \(x\) and numerical radii \(a,b\), say that the ball is
\emph{interior to \(H\)} if
\begin{equation}
\label{hc:interior}
 a+\delta_i+\Delta_i\le\|z-x\|_1
 \le b-\delta_i-\Delta_i.
\end{equation}
An interior ball is contained in \(H\), by the triangle inequality.
It is nonempty, so disjointness of \(\mathcal Z\) makes its containing
object unique.  Moreover, this object cannot be fresh: an unprocessed
ball and a fresh ball belong to the same family, whose balls are
disjoint.

While some unprocessed bag has an interior ball, choose such a bag
\(D=\{B_Q(z_j,\delta_i):1\le j\le n\}\).  Let \(H\) be the old
object containing an interior ball of \(D\), and set
\[
 L=\min_j\|z_j-x\|_1,
 \qquad R=\max_j\|z_j-x\|_1.
\]
The function \(z\mapsto\|z-x\|_1\) is 1-Lipschitz.  The diameter
bound on the centers, applied first to the interior center and then
to the extremal pair, gives
\begin{equation}
\label{hc:radial-span}
 a+\delta_i\le L\le R\le b-\delta_i,
 \qquad R-L\le\Delta_i.
\end{equation}
Thus every ball of \(D\) lies in \(H\).
Replace \(H\) by the two old shells
\begin{equation}
\label{hc:split}
 H_-=B_Q(x;a,L-\delta_i),
 \qquad H_+=B_Q(x;R+\delta_i,b),
\end{equation}
discarding any zero-width shell, and insert all \(n\) balls of \(D\)
as fresh balls.  Mark \(D\) processed.

Here are the geometric invariant checks.  The shells in
\eqref{hc:split} are subsets of \(H\) and are disjoint from each other.
For \(y\in B_Q(z_j,\delta_i)\), strictness of the ordinary ball gives
\[
 L-\delta_i<\|y-x\|_1<R+\delta_i.
\]
It therefore lies in neither shell.  The inserted balls are mutually
disjoint by the family assumption; all replacements are subsets of
\(H\), so they also avoid every other live object.
For \(H_-\), an old inner witness remains an inner witness, and a
center achieving \(L\) is an outer witness.  For \(H_+\), a center
achieving \(R\) is an inner witness, and an old outer witness remains
an outer witness.  Every inserted ordinary ball has its own center
as inner witness and the center of another ball of \(D\) as outer
witness; their distance is at least \(2\delta_i\).
All centers of new objects are either \(x\) or a \(z_j\), hence are
embedded vertices.  This proves (i) and (ii), and also proves that
fresh objects have not been altered.

For the potential calculation, write \(w=b-a\) and
\(t=[w-6\Delta_i]_+\), the old capacity being replaced.  The new
ordinary balls have total capacity
\(n(\delta_i-6\Delta_{i+1})\ge n\delta_i/2\).
If \(t=0\), they alone pay the new requirement \(n\delta_i/4\),
and no previous capacity is lost.  If \(t>0\), then
\(t=w-6\Delta_i\).  The widths of the two shells sum to
\(w-(R-L)-2\delta_i\), including zero-width shells formally, so
\begin{align}
\label{hc:split-capacity}
 &p_i(H_-)+p_i(H_+)+n(\delta_i-6\Delta_{i+1})\notag\\
 &\qquad\ge w-(R-L)-2\delta_i-12\Delta_i
              +\frac{n\delta_i}{2}\notag\\
 &\qquad\ge t+\frac{n\delta_i}{2}-8\Delta_i
 \ge t+\frac{n\delta_i}{4}.
\end{align}
The first bound uses \([s]_+\ge s\), also when a shell is discarded.
The second uses \(R-L\le\Delta_i\) and
\(2\delta_i\le\Delta_i\).  The last uses
\(n\delta_i\ge \beta\Delta_i\) and \(\beta\ge36\ge32\).
Unchanged objects keep their capacities.  Thus (iii) also survives,
with the new bag included on its right-hand side.  Each such operation
processes one new bag, so the loop ends after finitely many steps.

\paragraph{Inserting the remaining bags and changing scale.}
At the end of the loop, let \(\mathcal R_i\) be the unprocessed bags.
None of their balls satisfies \eqref{hc:interior} for any live object.
Set
\(s_i=2\delta_i+\Delta_i\).
Retain every fresh ball.  For each old object with numerical radii
\(a,b\), replace it by
\begin{equation}
\label{hc:shrink}
 H^{\mathrm{sh}}=B_Q(x;a+s_i,b-s_i)
 \quad\text{if }b-a>2s_i,
\end{equation}
and otherwise discard it.  Insert every ball of every bag in
\(\mathcal R_i\), and mark those bags processed.  We now regard all
surviving objects as old objects at the beginning of phase \(i+1\).
After phase \(q\), this is just a final accounting state with
\(\Delta_{q+1}=0\) and capacity equal to width.

Shrinking preserves pairwise disjointness among old objects and
preserves their witnesses: increasing the inner radius and decreasing
the outer radius only weakens the witness conditions.  A remaining
ball \(B_Q(z,\delta_i)\) is disjoint from each retained fresh ball,
because both belong to \(\mathcal F_i\).  For an old object, failure
of \eqref{hc:interior} says that at least one of
\[
 \|z-x\|_1<a+\delta_i+\Delta_i,
 \qquad
 \|z-x\|_1>b-\delta_i-\Delta_i
\]
holds.  In the first case every point of \(B_Q(z,\delta_i)\) has
radial distance less than \(a+s_i\); in the second every such point
has radial distance greater than \(b-s_i\).  In either case this ball
avoids \(H^{\mathrm{sh}}\).  All newly inserted balls are mutually
disjoint by the family assumption, and their witnesses are again
their own centers and another center in the same bag.  This verifies
(i) and (ii) after the transition, including the center condition.

It remains to verify (iii) across the change in its definition of
capacity.  A retained fresh ball has capacity
\(\delta_i-6\Delta_{i+1}\) both before and after the transition.
For a surviving old object, the new capacity is
\([
w-2s_i-6\Delta_{i+1}]_+\).  By \eqref{hc:scales},
\begin{equation}
\label{hc:reserve}
 2s_i+6\Delta_{i+1}
 =2\Delta_i+4\delta_i+6\Delta_{i+1}
 \le2\Delta_i+\frac92\delta_i
 \le\frac{17}{4}\Delta_i
 \le6\Delta_i.
\end{equation}
Thus its new capacity is at least \([w-6\Delta_i]_+\), its old
capacity.  If an old object is discarded, then
\(w\le2s_i\le4\Delta_i\), so its old capacity was zero.
Finally, every ball inserted from \(\mathcal R_i\) contributes
\(\delta_i-6\Delta_{i+1}\ge\delta_i/2\) in the new state.  These
capacities pay at least \(|D|\delta_i/4\) for each newly processed
bag \(D\).  No previous capacity decreases, so
\eqref{hc:potential-invariant} holds at the start of the next phase.

After phase \(q\), every bag in the chosen initial sequence has been
processed and the capacity of each live object equals its width.
The final collection \(\mathcal Z_{\mathrm{fin}}\) is disjoint and
has both witnesses for every object.  The potential invariant and
Lemma~\ref{hc:length} give
\[
 \frac14\sum_{i=1}^q\delta_i\sum_{D\in\mathcal F_i}|D|
 \le\sum_{H\in\mathcal Z_{\mathrm{fin}}}w(H)
 \le\frac1k\sum_{e\in E}\|X_e\|_1.
\]
For a countable sequence, its nonnegative sum is the supremum of
these finite initial sums.  Taking that supremum proves
\eqref{hc:conclusion}.
\end{proof}
